
\chapter{M002 --- Primitive Surface Construction}

\section{Stage 3 --- Implementation Review (Phase 5)}

\subsection{Reviewed Transformation}

Stacking-vector shear minimization while preserving the primitive-surface
construction constraints.

\subsection{Mathematical Input}

\begin{itemize}
\item Certified primitive in-plane basis.
\item Certified Bézout witness.
\item Stacking vector satisfying the surface-frame constraints.
\end{itemize}

\subsection{Mathematical Output}

A shear-minimized stacking vector representing the same crystallographic
surface and preserving the primitive construction certificate.

\subsection{Algorithmic Role}

This phase improves the numerical and geometric quality of the slab
representation without changing the represented primitive surface. The
implementation seeks a preferred representative within an equivalence
class rather than constructing a new scientific object.

\subsection{Classification}

\begin{itemize}
\item Gauge transformation.
\item Canonical representative selection.
\item Representation improvement.
\end{itemize}

\subsection{Preserved Invariants}

\begin{itemize}
\item Primitive kernel certificate.
\item Bézout constraint.
\item Surface orientation.
\item Primitive surface lattice.
\end{itemize}

\subsection{Open Questions for Stage 4}

\begin{enumerate}
\item Is the shear minimization mathematically equivalent to a gauge transformation?
\item Does every update preserve the Bézout witness exactly?
\item Is the selected representative optimal under the documented metric?
\end{enumerate}
